% 03-simulations.tex — État de l'art 2/5 : effets des simulations sur
% l'apprentissage, guidage, charge cognitive, représentations multiples,
% outils existants.

\chapter{Simulations interactives et apprentissage}\label{chap:simulations}

Ce chapitre présente ce que la recherche montre de l'apprentissage avec des
simulations, les conditions de leur efficacité, puis les outils qui existent
déjà.

\section{Effets des simulations sur l'apprentissage}

Les simulations aident à apprendre, à condition d'être bien employées. La
revue de Rutten, van Joolingen et van der Veen, qui couvre une décennie
d'études en sciences, conclut que les simulations améliorent l'enseignement
classique quand elles le complètent, en particulier pour préparer ou
accompagner une activité pratique \cite{rutten2012}. Plus largement, la
méta-analyse de Freeman et al.~montre que l'apprentissage actif, dont les
simulations sont une forme possible, améliore les résultats aux examens et
réduit l'échec en sciences et en ingénierie \cite{freeman2014}.

Le projet PhET, à l'université du Colorado, est le plus connu du domaine.
Ses simulations ont été conçues puis testées en entretiens avec des
étudiants. Les auteurs en tirent quelques principes : une réponse immédiate à
chaque action, plusieurs représentations du même phénomène, et une
exploration laissée ouverte mais orientée par l'interface
\cite{wieman2008phet, adams2008phet}. Finkelstein et al.~ont même montré que,
pour les circuits électriques, des étudiants qui avaient travaillé avec une
simulation réussissaient mieux à monter un vrai circuit que ceux qui avaient
manipulé du vrai matériel \cite{finkelstein2005}. Une simulation peut en effet
montrer ce qu'un laboratoire cache : le courant, les forces, le champ.

\section{Le guidage}

Laisser l'étudiant découvrir seul ne suffit pas. De~Jong et van~Joolingen ont
passé en revue l'apprentissage par découverte avec des simulations : les
étudiants formulent mal leurs hypothèses, changent plusieurs variables à la
fois et tirent des conclusions de trop peu d'essais \cite{dejong1998}. La
découverte devient efficace quand l'environnement guide : hypothèses à
tester proposées, une seule variable modifiée à la fois, un objectif à
atteindre.

Pour une animation, cela se traduit par un verdict clair (le wagon fait le
tour, ou décolle à tel angle) et par des questions ou des défis ciblés.
Afficher les hypothèses du \emph{modèle} (frottements négligés, objet
assimilé à un point, corde inextensible) est une exigence voisine mais
distincte : l'étudiant doit savoir ce que la simulation idéalise pour juger
ce qu'elle montre.

\section{Charge cognitive et principes multimédias}

Une simulation riche peut aussi gêner. La théorie de la charge cognitive part
d'un fait : la mémoire de travail est limitée, et tout ce qui l'occupe sans
servir l'apprentissage le ralentit \cite{sweller1988}. Mayer en a tiré des
principes de conception vérifiés expérimentalement \cite{mayer2009}. Trois
concernent directement une page de simulation :
\begin{itemize}
  \item \emph{contiguïté} : un texte et la figure qu'il commente doivent être
        proches dans l'espace et dans le temps, pour que l'étudiant n'ait pas
        à garder l'un en mémoire en cherchant l'autre ;
  \item \emph{signalement} : des repères visuels constants indiquent ce qui
        compte, par exemple une couleur fixe par type de force ;
  \item \emph{cohérence} : ce qui n'apporte rien à la compréhension (décor,
        effets) doit être retiré.
\end{itemize}
Ces principes justifient plusieurs de nos choix : tout tient sur un écran
quand c'est possible, chaque force a sa couleur dans tous les exercices
(poids en rouge, réaction normale en vert, tension en violet…), et les
formules sont placées juste sous la scène.

\section{Représentations multiples : schéma, formules, graphes}

En physique, le même phénomène se lit sous plusieurs formes : la scène, le
diagramme du corps libre, l'équation, le graphe en fonction du temps.
Ainsworth montre que ces représentations multiples aident quand elles se
complètent ou s'éclairent l'une l'autre, mais que passer de l'une à l'autre
est une tâche en soi, que les débutants réussissent mal sans aide
\cite{ainsworth2006deft}.

Une animation peut aider à faire ce lien : toutes les vues avancent au même
instant, et la formule affiche les valeurs du moment présent. C'est l'idée de
la formule écrite trois fois (symboles, valeurs substituées, résultat) : elle
relie le calcul de l'exercice à ce que montre la scène. Ce lien automatique a
un revers : il peut dispenser l'étudiant de faire la traduction lui-même
\cite{ainsworth2006deft}. D'où l'intérêt d'activités où il doit prédire avant
de voir, comme les défis du mode jeu (chapitre~\ref{chap:ludification}).

\section{Outils existants}

De nombreux outils proposent déjà des simulations de mécanique :
\begin{itemize}
  \item \textbf{PhET} : simulations en HTML5, gratuites et traduites,
        conçues pour un public large, du secondaire à l'université
        \cite{phet_site} ;
  \item \textbf{Easy Java/JavaScript Simulations (EjsS)} : un environnement
        qui permet à un enseignant de construire une simulation à partir de
        ses équations, sans écrire toute l'interface \cite{esquembre2004ejs} ;
  \item \textbf{myPhysicsLab} : des simulations de mécanique à code ouvert,
        avec intégration numérique explicite et graphes \cite{myphysicslab} ;
  \item \textbf{GeoGebra} : un logiciel de géométrie dynamique dans lequel un
        enseignant construit lui-même ses animations \cite{geogebra} ;
  \item \textbf{les animations du cours LEPL1201} : \enquote{Vitesse et
        accélération} et \enquote{Comète de Halley}, écrites en JavaScript
        pour le cours \cite{zoulab_epl1201}.
\end{itemize}

Le tableau~\ref{tab:outils} compare ces familles d'outils sur les besoins
propres à un cours. Il ne juge pas leur qualité générale, qui est bonne, mais
leur adéquation à LEPL1201.

\begin{table}[ht]
  \centering
  \small
  \caption{Adéquation des outils existants à un cours précis. \enquote{À
    construire} : possible, mais l'enseignant doit le programmer lui-même.}
  \label{tab:outils}
  \begin{tabularx}{\linewidth}{@{}>{\raggedright\arraybackslash}p{4.1cm} *4{>{\centering\arraybackslash}X}@{}}
    \toprule
    & Simulations générales & Outils de construction & Animations du cours
    & Notre site \\
    & (PhET, myPhysicsLab) & (EjsS, GeoGebra) & (LEPL1201) & \\
    \midrule
    Notations du formulaire        & non      & à construire & oui & oui \\
    Calcul avec valeurs substituées & non     & à construire & non & oui \\
    Forces et diagramme du corps libre & selon la simulation & à construire
                                   & \(\vec v\), \(\vec a\) seulement & oui \\
    Nombre de situations           & grand    & --           & 2   & 5 \\
    Effort pour un nouvel exercice & --       & élevé        & élevé & à mesurer \\
    \bottomrule
  \end{tabularx}
\end{table}

À part les animations du cours, aucun de ces outils n'est écrit pour un cours
précis. Les simulations générales n'emploient pas les notations du
formulaire, n'affichent pas le calcul de l'exercice avec ses valeurs, et sont
difficiles à modifier par un enseignant. Les outils de construction peuvent
être adaptés, mais au prix du temps de développement que nous cherchons
justement à réduire. Notre travail ne cherche donc pas à refaire PhET : il
vise des animations écrites pour un cours, dans ses notations, qu'un
enseignant peut produire et modifier lui-même. La dernière ligne du tableau
est une hypothèse, que la partie Résultats doit vérifier.
